\documentclass[a4paper,11pt]{article}
\usepackage[T1]{fontenc}
\usepackage[polish]{babel}
\usepackage[utf8]{inputenc}
\usepackage{geometry}
\usepackage{listings}
\usepackage{xcolor}
\usepackage{parskip}
\usepackage{titlesec}
\usepackage{hyperref}

\geometry{margin=2.5cm}

\definecolor{codebg}{RGB}{245,245,245}
\definecolor{kw}{RGB}{0,100,180}
\definecolor{str}{RGB}{160,32,32}
\definecolor{comment}{RGB}{100,120,100}

\lstdefinelanguage{Kotlin}{
  keywords={fun, val, listOf, data, class, object, null, true, false, override, private},
  keywordstyle=\color{kw}\bfseries,
  stringstyle=\color{str},
  commentstyle=\color{comment}\itshape,
  morestring=[b]",
  morecomment=[l]{//},
  basicstyle=\ttfamily\small,
  backgroundcolor=\color{codebg},
  frame=single,
  framerule=0.4pt,
  rulecolor=\color{gray!40},
  xleftmargin=8pt,
  xrightmargin=8pt,
  breaklines=true,
  showstringspaces=false,
}

\lstset{language=Kotlin}

\titleformat{\section}{\large\bfseries}{}{0em}{}[\titlerule]
\titleformat{\subsection}{\normalsize\bfseries}{}{0em}{}

\title{\textbf{RUTEdu}}
\author{}
\date{}

\begin{document}
\maketitle
\tableofcontents
\vspace{1em}

% ─────────────────────────────────────────────────────────────
\section{Struktura projektu}

Projekt to \textbf{Kotlin Multiplatform} (Android + iOS) z Jetpack Compose.
Cały wspólny kod leży w \texttt{composeApp/src/commonMain/kotlin/}.

\begin{center}
\begin{tabular}{ll}
\texttt{data/SubjectRepository.kt} & lista przedmiotów, tematów i lekcji \\
\texttt{data/QuestionBank.kt}      & statyczne pytania (mat, geo, chemia\_3) \\
\texttt{data/ChemistryQuestionGenerator.kt} & generowane pytania chemiczne \\
\texttt{models/Question.kt}        & typy pytań (sealed class) \\
\texttt{screens/}                  & komponenty UI dla każdego typu pytania \\
\end{tabular}
\end{center}

% ─────────────────────────────────────────────────────────────
\section{Jak dodawać nowe zadania}

\subsection{Krok 1 - zarejestruj lekcję w SubjectRepository.kt}

Otwórz \texttt{SubjectRepository.kt} i znajdź właściwy temat (\texttt{Topic}).
Dodaj nowy wpis \texttt{Lesson} do jego listy \texttt{lessons}:

\begin{lstlisting}
Lesson(
    id          = "mat_2_4",          // unikalny string, uzywany jako klucz
    name        = "Ulamki dziesietne",
    description = "Dodawanie i odejmowanie",
    progress    = 0f,
    isLocked    = false,
    icon        = Icons.Default.Calculate,
    lessonType  = LessonType.DEFAULT
)
\end{lstlisting}

\textbf{Zasady ID:} format \texttt{<przedmiot>\_<temat>\_<numer>}, np.\ \texttt{mat\_2\_4}, \texttt{geo\_3\_1}, \texttt{chemia\_2\_3}.

\subsection{Krok 2a - dodaj pytania statyczne (QuestionBank.kt)}

Dla matematyki, geografii i \texttt{chemia\_3\_x} pytania trafiają do \texttt{QuestionBank.kt}.

\begin{enumerate}
  \item Stwórz prywatną listę pytań:
\begin{lstlisting}
private val mat_2_4: List<Question> = listOf(
    TypeAnswer(
        id            = 0,
        prompt        = "Oblicz: 1.5 + 2.3 = ?",
        correctAnswer = 38,   // wynik * 10 jesli jednostka "dziesiete"
        unit          = "",
        hint          = Hint("Dodaj oddzielnie czesci cale i dziesietne.")
    ),
    SelectFromList(
        id             = 1,
        prompt         = "Ktory ulamek jest wiekszy?",
        options        = listOf("0.7", "0.3", "0.5"),
        correctIndices = setOf(0),
        hint           = Hint("Porownaj pierwsza cyfre po przecinku.")
    ),
    // ... minimum 8-10 pytan
)
\end{lstlisting}

  \item Zarejestruj liste w mapie \texttt{banks} (na dole pliku):
\begin{lstlisting}
private val banks: Map<String, List<Question>> = mapOf(
    // ... istniejace lekcje ...
    "mat_2_4" to mat_2_4,
)
\end{lstlisting}
\end{enumerate}

\subsection{Krok 2b - wygenerowane pytania chemiczne (ChemistryQuestionGenerator.kt)}

Jesli ID lekcji zaczyna sie od \texttt{chemia\_} (i nie jest \texttt{chemia\_3\_x}),
\texttt{QuestionBank} automatycznie przekierowuje do generatora.

\begin{enumerate}
  \item Dodaj case w funkcji \texttt{generateFor()}:
\begin{lstlisting}
fun generateFor(lessonId: String, seed: Long = ...): List<Question> =
    when (lessonId) {
        // istniejace przypadki...
        "chemia_2_3" -> chemia_2_3(seed)
        else         -> emptyList()
    }
\end{lstlisting}

  \item Zaimplementuj funkcje generujaca:
\begin{lstlisting}
private fun chemia_2_3(seed: Long): List<Question> {
    val rng = Random(seed)
    val wybrane = ELEMENTS.shuffled(rng).take(10)
    return wybrane.mapIndexed { i, el ->
        ElementCardQuiz(
            id           = i,
            prompt       = "Jaka jest masa atomowa tego pierwiastka?",
            atomicNumber = el.atomicNumber,
            options      = generateMassOptions(el, rng),
            correctIndex = 0,
            hint         = Hint("Masa atomowa to srednia wazonych izotopow.")
        )
    }
}
\end{lstlisting}

  Dostepne dane: \texttt{ELEMENTS} (lista wszystkich 118), \texttt{elementByNumber}, \texttt{shellConfigByNumber}.
\end{enumerate}

% ─────────────────────────────────────────────────────────────
\section{Dostepne typy pytan}

\begin{center}
\renewcommand{\arraystretch}{1.25}
\begin{tabular}{lp{8cm}}
\textbf{Typ} & \textbf{Opis} \\
\hline
\texttt{FindAnswer}         & Uzytkownik wpisuje wynik dzialania liczbowego \\
\texttt{FindOperator}       & Uzytkownik przeciaga operator (+, -, ×, ÷) \\
\texttt{SelectFromList}     & Wybor jednej lub kilku opcji z listy (radio/checkbox) \\
\texttt{TypeAnswer}         & Wpisanie liczby; opcjonalny diagram trojkata \\
\texttt{MapQuiz}            & Klikniecie kraju na mapie (MapLibre) \\
\texttt{PeriodicTableQuiz}  & Umiesc brakujacy pierwiastek w tablicy \\
\texttt{PeriodicTableByShell} & Znajdz pierwiastek po konfiguracji powlokowej \\
\texttt{PeriodicTableByName}  & Znajdz pierwiastek po nazwie polskiej \\
\texttt{ElementCardQuiz}    & Karta pierwiastka, wybor 1 z 4 odpowiedzi \\
\texttt{EquationBalance}    & Uzupelnienie wspolczynnikow stechiometrycznych \\
\texttt{GraphTypeAnswer}    & Wykres + wpisanie wyniku liczbowego \\
\texttt{GraphSelectFromList}& Wykres + wybor z listy \\
\end{tabular}
\end{center}

\subsection{Przyklad: Hint}

Kazde pytanie przyjmuje opcjonalny \texttt{Hint} wyswietlany w dolnym panelu:

\begin{lstlisting}
Hint(
    mainText     = "Wzor: V = a * b * h",
    boldPart     = "V = a * b * h",         // fragment pogrubiany
    sectionTitle = "Zapamietaj",
    items        = listOf("a, b, h to krawedzie"),
    steps        = listOf("Zmierz a", "Zmierz b", "Zmierz h", "Pomnoz")
)
\end{lstlisting}

% ─────────────────────────────────────────────────────────────
\section{Przydatne komendy}

\begin{lstlisting}[language={}]
# Build Android
./gradlew :composeApp:assembleDebug

# Uruchom na podlaczonym urzadzeniu
./gradlew :composeApp:installDebug

# Sprawdz bledy kompilacji
./gradlew :composeApp:compileKotlinAndroid
\end{lstlisting}

\end{document}
